%% Owner: Víctor Martín Castilla
\section{Infrastructure}
\label{sec:infrastructure}

\todo{Rewrite for the cluster the team uses: three Azure VMs built by
Mikael. Document where they run, meaning subscription type, region, VM
size and operating system. Document the network: address space, which
VMs have public IPs, and the firewall rules. Then how the VMs were
created and prepared, ideally as code; access control and SSH
hardening, with evidence; the options the team considered for hosting;
and the cost. Text about an earlier cluster, built and deleted on
17 October, is in the git history if anything is worth reusing.}

\subsection{Requirements}
\label{sec:infra-requirements}

The statement asks for a Swarm of at least three nodes and leaves the
hosting open: home hardware, virtual machines on a laptop, or a public
cloud \parencite{llusa2026project}. The choice had to meet five
constraints:
\begin{itemize}
  \item three machines that stay on, so the cluster works whenever it is
        checked;
  \item a public address for the front proxy, so Let's Encrypt can
        validate the domain and users can reach the services;
  \item a private network between the nodes for Swarm and NFS traffic;
  \item no cost for the team;
  \item ready within the two days left before the deadline.
\end{itemize}

\subsection{Why three nodes}
\label{sec:infra-three}

Swarm managers keep the cluster state with the Raft consensus algorithm
\parencite{ongaro2014raft,llusa2026orchestration}. A change is committed
only when a majority of managers agree, so a cluster of $N$ managers
survives $\lfloor (N-1)/2 \rfloor$ manager failures
\parencite{docker-raft}. \Cref{tab:raft} shows why two managers are no
better than one, and why three is the smallest cluster that survives a
failure. All three nodes should therefore be managers, and they also run
workloads.

\begin{table}[htbp]
  \centering
  \caption{Manager failures a Swarm survives while keeping a majority.}
  \label{tab:raft}
  \begin{tabular}{ccc}
    \toprule
    Managers & Majority & Failures survived \\
    \midrule
    1 & 1 & 0 \\
    2 & 2 & 0 \\
    3 & 2 & 1 \\
    5 & 3 & 2 \\
    \bottomrule
  \end{tabular}
\end{table}
